\documentclass[17pt, orivec]{extarticle} % 12pt, 14pt, 17pt, 20pt
\usepackage{amsmath}
\usepackage{amssymb}    % for \rightsquigarrow
\usepackage{wasysym}	% for frown face
\usepackage{mathrsfs} 	% for \mathscr
\usepackage{stmaryrd}
\usepackage[most]{tcolorbox}
\usepackage{ulem}
\usepackage{tikz-cd}		% commutative diagrams
\usepackage{tikz}
\usepackage{amsthm}
\usepackage{enumitem}	% for \itemize custom labels
\usepackage{turnstile}	% longer turnstiles
\usepackage[sf,bf,big,raggedright,compact]{titlesec}	% change section color to blue
% \usepackage[backend=biber,bibstyle=authoryear,citestyle=../authoryearbrack]{biblatex}
% \bibliography{../AGI-book}

\newtheorem{theorem}{Theorem}

\ifdefined\chinchin
	\newcommand{\cc}[2]{#1}
	\usepackage[CJKspace]{xeCJK}
	%\setCJKmainfont[BoldFont=SimHei,ItalicFont=AR PL KaitiM GB]{Alibaba PuHuiTi}
	\setCJKmainfont[BoldFont=Alibaba-PuHuiTi-Regular.ttf, ItalicFont=gkai00mp.ttf]{Alibaba-PuHuiTi-Light.ttf}
	% \setmainfont[ItalicFont=Latin Modern Roman Slanted]{Alibaba Sans:style=Light}
\else
	\newcommand{\cc}[2]{#2}
	% \setmainfont[ItalicFont=Latin Modern Roman Slanted]{Alibaba Sans:style=Light}
	%\renewcommand{\baselinestretch}{0.8} 
\fi

%\ifdefined\chinchin
%\newcommand{\cc}[2]{#1}
%\usepackage[CJKspace]{xeCJK}
%\setCJKmainfont[BoldFont=SimHei,ItalicFont=AR PL KaitiM GB]{SimSun}
%\else
%\newcommand{\cc}[2]{#2}
%\fi

\setlength{\headheight}{0cm}
\setlength{\hoffset}{0cm}
\setlength{\topmargin}{-2cm}
\setlength{\oddsidemargin}{-2cm}
\setlength{\evensidemargin}{-2cm}
\setlength{\textwidth}{20cm}
\setlength{\textheight}{26cm}
\setlength{\headsep}{0cm}
\setlength{\topskip}{0cm}
\setlength{\footskip}{0.9cm}  % between bottom of page and page number
\setlength{\floatsep}{0cm}
\setlength{\textfloatsep}{0.6cm}
\setlength{\intextsep}{0.5cm}
\setlength{\parindent}{0em}   % em = width of capital M

% Fix spilling of titles in bibliography:
%\DeclareFieldFormat*{title}{#1}
%
%\DeclareFieldFormat*{titlecase}{%
%    \ifdef{\currentfield}
%      {\ifcurrentfield{title}
%         {\usefield{\uline}{\currentfield}}%
%         {#1}}
%      {#1}}

\setcounter{secnumdepth}{2}		% no section numbers

\titleformat{\section}[hang]{\bfseries\Large\color{blue}}{\thesection \hspace{20pt}}{0pt}{}
\titleformat{\subsection}[hang]{\bfseries\large\color{blue}}{\thesubsection \hspace{10pt}}{0pt}{}
\titleformat{\subsubsection}[hang]{\bfseries\color{blue}}{}{0pt}{}

\itemsep0em
\setlist[itemize]{noitemsep, topsep=-5pt, partopsep=-5pt}
\renewcommand{\labelitemi}{$\textbullet$}

\let\varzero\emptyset
\let\emptyset\varnothing		% round empty set 
\newcommand{\vect}[1]{\boldsymbol{#1}}
\newcommand{\book}[1]{$\NewSym[0.4]{../book-icon.png} \quad$ \parbox{0.9\textwidth}{\footnotesize #1}}
\newcommand{\code}[1]{{\footnotesize{\ttfamily #1}}}
\newcommand{\tab}{\hspace*{2cm}}
\newcommand{\powerset}{\raisebox{.15\baselineskip}{\Large\ensuremath{\wp}}}
\newcommand{\Chi}{\raisebox{2.5pt}{$\chi$}}
\newcommand*\KB{\vcenter{\hbox{\includegraphics{../KB-symbol.png}}}}
\newcommand*\NewSym[2][0.5]{\vcenter{\hbox{\includegraphics[scale=#1]{#2}}}}
\newcommand*\sigmoid{\vcenter{\hbox{\includegraphics{../sigmoid3.png}}}}
\newcommand{\smbox}[1]{\boxed{\footnotesize{\mbox{#1}}}}

\newcommand{\tikzmark}[1]{\tikz[overlay,remember picture] \node (#1) {};}

\newcommand{\Dfrac}[2]{%
\ooalign{%
      $\genfrac{}{}{2.9pt}0{\hphantom{#1}}{\hphantom{#2}}$\cr%
      $\color{white}\genfrac{}{}{1.5pt}0{\hphantom{#1}}{\hphantom{#2}}$\cr%
      $\color{white}\genfrac{}{}{1pt}0{\color{black}#1}{\color{black}#2}$}}

% \renewcommand{\thefootnote}{\fnsymbol{footnote}}
\usepackage{perpage}
\MakePerPage{footnote}
\renewcommand{\thefootnote}{\ifcase\value{footnote}\or{*}
	\or{$\dagger$}\or{**}\or{$\ddagger$}\fi}
\interfootnotelinepenalty=10000

\setlength{\parindent}{0pt}
\setlength{\parskip}{1.8ex plus0.8ex minus0.8ex}

% \usepackage[no-math]{fontspec}
% \setmainfont[BoldFont=Alibaba_Sans_Regular.otf,ItalicFont=Alibaba_Sans_Light_Italic.otf]{Alibaba_Sans_Light.otf}

\usepackage[backend=biber]{biblatex}
\bibliography{../AGI-book}

\usepackage[active,tightpage]{preview}		% for continuous page(s)
\renewcommand{\PreviewBorder}{0.5cm}
\renewcommand{\thempfootnote}{\arabic{mpfootnote}}

\usepackage[absolute,overlay]{textpos}		% for page number on upper left corner

\usepackage{color}
% \usepackage{mathtools}
\usepackage[hyperfootnotes=false]{hyperref}

% \usepackage[backend=biber,style=numeric]{biblatex}
% \bibliography{../AGI-book}
% \renewcommand*{\bibfont}{\footnotesize}

\usetikzlibrary{shapes}
% \usepackage[export]{adjustbox}	% ??
\usepackage{verbatim} % for comments
% \usepackage{newtxtext,newtxmath}	% Times New Roman font

\titleformat{\section}[hang]{\bfseries\large\color{cyan}}{}{0pt}{}
% \numberwithin{equation}{subsection}

\newcommand{\underdash}[1]{%
	\tikz[baseline=(toUnderline.base)]{
		\node[inner sep=1pt,outer sep=10pt] (toUnderline) {#1};
		\draw[dashed] ([yshift=-0pt]toUnderline.south west) -- ([yshift=-0pt]toUnderline.south east);
	}%
}%

\newcommand\reduline{\bgroup\markoverwith{\textcolor{red}{\rule[-0.5ex]{2pt}{0.4pt}}}\ULon}

%\DeclareSymbolFont{symbolsC}{U}{txsyc}{m}{n}
%\DeclareMathSymbol{\strictif}{\mathrel}{symbolsC}{74}
%\DeclareSymbolFont{AMSb}{U}{msb}{m}{n}
%\DeclareSymbolFontAlphabet{\mathbb}{AMSb}
%\setmathfont{lmroman17-regular.otf}
\DeclareMathOperator*{\argmin}{arg\,min}
\DeclareMathOperator*{\argmax}{arg\,max}

% \usepackage[most]{tcolorbox}
%\tcbset{on line,
%	boxsep=4pt, left=0pt,right=0pt,top=0pt,bottom=0pt,
%	colframe=red,colback=pink,
%	highlight math style={enhanced}
%}
%\newcommand{\atom}{\vcenter{\hbox{\tcbox{....}}}}

\let\oldtextbf\textbf
\renewcommand{\textbf}[1]{\textcolor{cyan}{\oldtextbf{#1}}}

\newcommand{\logic}[1]{{\color{violet}{\textit{#1}}}}
\newcommand{\underconst}{\includegraphics[scale=0.5]{../2020/UnderConst.png}}
\newcommand{\KBsymbol}{\vcenter{\hbox{\includegraphics[scale=1]{../KB-symbol.png}}}}
\newcommand{\token}{\vcenter{\hbox{\includegraphics[scale=1]{token.png}}}}
\newcommand{\proposition}{\vcenter{\hbox{\includegraphics[scale=0.8]{proposition.png}}}}

\begin{document}

\begin{preview}

\cc{
\title{\vspace{-1cm} \bfseries\color{cyan}{\Large Solidi 加密貨幣 - 白皮書}}
}{
\title{\vspace{-1cm} \bfseries\color{cyan}{\LARGE Solidi Coin - White Paper}}
}

% \cc{
% \author{YKY (甄景賢)}}{
% \author{YKY (Yan King Yin)}
% }

\date{\vspace{-1.5cm}\today} % Date, can be changed to a custom date

\cc{
\centerline{\includegraphics[scale=0.6]{凝聚力-logo-0.png}}
}{
\centerline{\includegraphics[scale=0.3]{Solidarity-logo-0.png}}
}

\maketitle

\setcounter{section}{-1}
\newcounter{mypage}
\setcounter{mypage}{0}

% (1) Circled page number on upper left corner
\begin{textblock*}{5cm}(2.1cm,2.3cm) % {block width} (coords)
{\color{red}{\large \textcircled{\small \themypage}}}
\addtocounter{mypage}{1}
\end{textblock*}

\begin{minipage}{\textwidth}
\setlength{\parskip}{0.4\baselineskip}

\section{\cc{Solidi 的貨幣經濟學\\}{Economics of the Solidi Coin \\}}

\begin{itemize}
	\item Solidi 幣的特點 是以\textbf{投票}形式 賺取的，例如 在新聞平台上發評論，被點贊即能獲得 貨幣獎勵。幣值也稱爲 \textbf{聲譽}，也等於 \textbf{投票權} (下述)。

\cc{
	\item \textbf{指數式 兌換：}  用戶可以花錢買 聲譽/投票權，但所需的價錢 會隨着聲譽 \textbf{指數式}增加。 換句話說，即使最有錢的個人或基金，都無法用錢換取極大的聲譽值，這保障了 聲譽不會純粹被最有錢的個人或財團主導。也可以說 Solidi 跟現時世界上的 money 是不同類的一種價值，它可以糾正 貧富懸殊的 財富分配，因爲現時世界的財富分配服從 log-normal distribution (對數-正態分佈)。 以下說「現金」泛指傳統金錢。
}{
	\item \textbf{Exponential Conversion.}  A member can spend money to ``pump up'' his reputation.  As reputation gets higher, it is increasingly harder to increase one unit of reputation, following an \textbf{exponential} curve: $e^{kR}$ is required to increase one unit of $R$, where $R$ is reputation, $k$ is a scaling constant, eg. $1/1000$.
}

\cc{
	\item \textbf{提取款項：}  用 Solidi 兌換現金，即上述的逆運算，則可以得到對數式的現金值。但這樣，當某會員得到了稍為可觀的聲譽，就很有可能 cash-out (withdraw)，獲得現金，而導致平台因缺乏現金而崩盤。 解決辦法是： cash-out 是根據 整體 聲譽的 \textbf{比例} 計算的。 當現金儲備 (即傳統 money) 很少的時候，cash-out 微不足道。 用一句口號概括就是：「指數式買入、線性賣出」
}{
	\item \textbf{Withdrawals} can be performed as the inverse conversion of the above.  But this will be problematic; see below.
}

\cc{
	\item \textbf{Solidi 的慷概性：}  賺取 聲譽 是很容易的，例如，每次 ``like'' 別人的評論，則評論者獲得 +1 聲譽值，而點贊者 扣除 $\frac{1}{100}$ 聲譽值。 因此人們不會吝嗇給別人點贊。
}{
	\item \textbf{Generosity of Likes.}  Every time a member likes a comment, the creator of the comment get +1 unit reputation, while the sender gets $\frac{1}{100}$ deduction in reputation, ie 0.01 unit.
}

\cc{}{
	\item Due to exponential conversion, when a member gets a considerable reputation, she may be tempted to cash-out (withdraw), causing the platform to run out of cash and collapse.  How to prevent this?
	}

\cc{}{
	\item A solution is:  withdrawals are calculated as a \textbf{proportion} of the reputation pool, instead of exponential conversion.  This means that when the total cash reserve of the platform is low, withdrawals would also be insignificant.
	}

\cc{
	\item 我們的核心價值是：
	\begin{equation} \nonumber
	\vcenter{\hbox{\includegraphics[scale=0.7]{DAO-核心價值.png}}}
	\end{equation}
	這些都是很經典的價值，在世界各地的文化裏都有其雛形，但在西方啓蒙時代系統性地提出，這並不表示盲目崇拜西方思想。 隨着科技進步，以上這些原則都是 \textbf{可驗證的} (verifiable)，例如 區塊鏈、開源軟件 等都是可驗證的。
}{
}

\cc{
	\item \textbf{投票權：} （這方面仍未想清楚） 一個科學家在他的研究領域是專家，但在別的領域不一定是專家，所以「聲譽」是有局部性的。 我們希望 Solidi 平台完全是由民主的方式運行的，也就是說，以 \textbf{加權投票} 決定平台的發展方向，而 權重 就是聲譽值。 這表示 投票權 必需以不同的 領域 劃分。
}{
}

\end{itemize}

\end{minipage}
\end{preview}

\begin{preview}
\begin{textblock*}{5cm}(2.1cm,2.3cm) % {block width} (coords)
	{\color{red}{\large \textcircled{\small \themypage}}}
	\addtocounter{mypage}{1}
\end{textblock*}

\begin{minipage}{\textwidth}
	\setlength{\parskip}{0.4\baselineskip}

\begin{itemize}
	\item

	\item

	\item
	\begin{itemize}
		\item
		\item
		\item

	\end{itemize}
\end{itemize}


\end{minipage}
\end{preview}

\begin{preview}
\begin{textblock*}{5cm}(2.1cm,2.3cm) % {block width} (coords)
{\color{red}{\large \textcircled{\small \themypage}}}
\addtocounter{mypage}{1}
\end{textblock*}

\begin{minipage}{\textwidth}
\setlength{\parskip}{0.4\baselineskip}


\end{minipage}
\end{preview}

\begin{preview}
\begin{textblock*}{5cm}(2.1cm,2.3cm) % {block width} (coords)
	{\color{red}{\large \textcircled{\small \themypage}}}
	\addtocounter{mypage}{1}
\end{textblock*}

\begin{minipage}{\textwidth}
	\setlength{\parskip}{0.4\baselineskip}

\section{DAO -- more even distribution of earnings}

\end{minipage}
\end{preview}

\end{document}
